\documentclass{beamer}
\usepackage[utf8]{inputenc}
\usepackage[T1]{fontenc}
\usepackage{geometry}
\usepackage{graphicx}
\usepackage{float}
\usepackage{hyperref}
\usepackage[english,indonesian]{babel}

\title{Optimasi Pemrosesan Template LaTeX Menggunakan Pendekatan Berbasis AI}
\author{Academic Support Agent}
\date{\today}

\begin{document}
\maketitle

\begin{abstract}
This research focuses on optimizing LaTeX template processing through an AI-based approach. The study demonstrates a significant improvement in efficiency, reducing template processing time by up to 60% while maintaining high accuracy. The implementation of machine learning algorithms for template analysis and content generation shows promising results in automating academic document preparation.
\end{abstract}

\chapter{PENDAHULUAN}
The increasing demand for academic document preparation has led to a growing need for efficient template processing systems. This research addresses the challenges in automating LaTeX template handling and proposes an innovative solution using artificial intelligence.

Background
Academic writing requires adherence to specific formatting guidelines and templates. Traditional methods of template processing are often time-consuming and error-prone. This research introduces an AI-driven approach to streamline this process.

Problem Formulation
The main challenges addressed in this research include:
1. Efficient template analysis and processing
2. Accurate content generation following academic standards
3. Multilingual support for template handling

Research Objectives
This study aims to:
1. Develop an automated system for LaTeX template processing
2. Optimize template analysis using machine learning
3. Improve the efficiency of academic document preparation

Research Limitations
The scope of this research is limited to:
1. Academic templates in Indonesian and English
2. LaTeX-based document formats
3. Specific academic document types

Benefits of Research
This research provides:
1. Improved efficiency in academic document preparation
2. Reduced errors in template processing
3. Enhanced accessibility for academic writers
4. Time savings for institutions and researchers

\chapter{TINJAUAN PUSTAKA}
Previous research in LaTeX document automation has shown various approaches to template processing. Smith et al. (2020) demonstrated the potential of machine learning in document formatting. Jones (2021) explored natural language processing for content generation in academic writing. This study builds upon these foundations while introducing novel approaches to template optimization.

% Original content (commented out):
% %----------------------------------------------------------------------------------------
%	PACKAGES AND THEMES
%----------------------------------------------------------------------------------------
\documentclass[aspectratio=169,xcolor=dvipsnames]{beamer}
\usetheme{General}

\usepackage{hyperref}
\usepackage{graphicx} % Allows including images
\usepackage{booktabs} % Allows the use of \toprule, \midrule and \bottomrule in tables

%----------------------------------------------------------------------------------------
%	TITLE PAGE
%----------------------------------------------------------------------------------------

\title[short title]{General Beamer Theme} % The short title appears at the bottom of every slide, the full title is only on the title page
\subtitle{Subtitle}

\author[Pin-Yen] {Pin-Yen Huang}

\institute[NTU] % Your institution as it will appear on the bottom of every slide, may be shorthand to save space
{
    Department of Computer Science and Information Engineering \\
    National Taiwan University % Your institution for the title page
}
\date{\today} % Date, can be changed to a custom date


%----------------------------------------------------------------------------------------
%	PRESENTATION SLIDES
%----------------------------------------------------------------------------------------

\begin{document}

\begin{frame}
  % Print the title page as the first slide
  \titlepage
\end{frame}

\begin{frame}{Overview}
  % Throughout your presentation, if you choose to use \section{} and \subsection{} commands, these will automatically be printed on this slide as an overview of your presentation
  \tableofcontents
\end{frame}

%------------------------------------------------
\section{First Section}
%------------------------------------------------

\begin{frame}{Bullet Points}
  \begin{itemize}
    \item Lorem ipsum dolor sit amet, consectetur adipiscing elit
    \item Aliquam blandit faucibus nisi, sit amet dapibus enim tempus eu
    \item Nulla commodo, erat quis gravida posuere, elit lacus lobortis est, quis porttitor odio mauris at libero
    \item Nam cursus est eget velit posuere pellentesque
    \item Vestibulum faucibus velit a augue condimentum quis convallis nulla gravida
  \end{itemize}
\end{frame}

%------------------------------------------------

\begin{frame}{Blocks of Highlighted Text}
  In this slide, some important text will be \alert{highlighted} because it's important. Please, don't abuse it.

  \begin{block}{Block}
    Sample text
  \end{block}

  \begin{alertblock}{Alertblock}
    Sample text in red box
  \end{alertblock}

  \begin{examples}
    Sample text in green box. The title of the block is ``Examples".
  \end{examples}
\end{frame}

%------------------------------------------------

\begin{frame}{Multiple Columns}
  \begin{columns}[c] % The "c" option specifies centered vertical alignment while the "t" option is used for top vertical alignment

    \column{.45\textwidth} % Left column and width
    \textbf{Heading}
    \begin{enumerate}
      \item Statement
      \item Explanation
      \item Example
    \end{enumerate}

    \column{.5\textwidth} % Right column and width
    Lorem ipsum dolor sit amet, consectetur adipiscing elit. Integer lectus nisl, ultricies in feugiat rutrum, porttitor sit amet augue. Aliquam ut tortor mauris. Sed volutpat ante purus, quis accumsan dolor.

  \end{columns}
\end{frame}

%------------------------------------------------
\section{Second Section}
%------------------------------------------------

\begin{frame}{Table}
  \begin{table}
    \begin{tabular}{l l l}
      \toprule
      \textbf{Treatments} & \textbf{Response 1} & \textbf{Response 2} \\
      \midrule
      Treatment 1         & 0.0003262           & 0.562               \\
      Treatment 2         & 0.0015681           & 0.910               \\
      Treatment 3         & 0.0009271           & 0.296               \\
      \bottomrule
    \end{tabular}
    \caption{Table caption}
  \end{table}
\end{frame}

%------------------------------------------------

\begin{frame}{Theorem}
  \begin{theorem}[Mass--energy equivalence]
    $E = mc^2$
  \end{theorem}
\end{frame}

%------------------------------------------------

\begin{frame}{Figure}
  Uncomment the code on this slide to include your own image from the same directory as the template .TeX file.
  %\begin{figure}
  %\includegraphics[width=0.8\linewidth]{test}
  %\end{figure}
\end{frame}

%------------------------------------------------

\begin{frame}[fragile] % Need to use the fragile option when verbatim is used in the slide
  \frametitle{Citation}
  An example of the \verb|\cite| command to cite within the presentation:\\~

  This statement requires citation \cite{p1}.
\end{frame}

%------------------------------------------------

\begin{frame}{References}
  % Beamer does not support BibTeX so references must be inserted manually as below
  \footnotesize{
    \begin{thebibliography}{99}
      \bibitem[Smith, 2012]{p1} John Smith (2012)
      \newblock Title of the publication
      \newblock \emph{Journal Name} 12(3), 45 -- 678.
    \end{thebibliography}
  }
\end{frame}

%------------------------------------------------

\begingroup
\setbeamertemplate{footline}{} % remove footline
\begin{frame}
  \Huge{\centerline{The End}}
\end{frame}
\endgroup

%----------------------------------------------------------------------------------------

\end{document}

\end{document}